\subsection{INVERSE IMAGE OF A UNIFORMITY; UNIFORM SUBSPACES}

Let X be a set, Y a uniform space, f a mapping of X into Y. The coarsest uniformity U on X for which f is uniformly continuous is called the inverse image under f of the uniform structure of Y. It follows from Proposition 4 of no. 3, and from the formulae which give the inverse image of an intersection, that the inverse images under \(g = f \times f\) of the entourages of Y form a fundamental system of entourages for U. The topology induced by U is the inverse image under f of the topology of Y (no. 3, Corollary to Proposition 4).

Remark. If \(f: \mathbf{X} \to \mathbf{Y}\) is surjective, then the entourages of \(\mathbf{Y}\) are the direct images under \(g\) of the entourages of \(\mathbf{X}\) .

A mapping \(f\) of a uniform space \(\mathbf{X}\) into a uniform space \(\mathbf{X}'\) is uniformly continuous if and only if the inverse image under \(f\) of the uniformity of \(\mathbf{X}'\) is coarser than the uniformity of \(\mathbf{X}\) .

Let A be a subset of a uniform space X. The uniformity induced on A by the uniformity of X is the inverse image of the latter under the canonical injection \(A \rightarrow X\) . By Proposition 4 of no. 3, this is equivalent to the following definition:

DEFINITION 3. Let A be a subset of a uniform space X. The uniformity on A whose set of entourages is the trace on A × A of the set of entourages of X is called the uniformity induced on A by the uniformity of X.

The topology induced by the uniformity induced on A is the same as the topology induced on A by the topology of X; the set A, together with the uniformity and the topology induced by those of X, is called a uniform subspace of X.

If A is a subset of a uniform space X and if \(f: X \to X'\) is a uniformly continuous mapping, then the restriction \(f|A\) is a uniformly continuous mapping of A into \(X'\) . If \(A' \subset X'\) is such that \(f(X) \subset A'\) , then the mapping of X into the uniform subspace \(A'\) of \(X'\) , having the same graph as f, is again uniformly continuous (no. 3, Proposition 4).

If \(B \subset A \subset X\) , then the uniform subspace B of X is identical with the uniform subspace B of the uniform subspace A of X (transitivity of induced uniform structures; no. 3, Proposition 5).

PROPOSITION 6. Let A be a dense subset of a uniform space X. Then the closures, in X × X, of the entourages of the uniform subspace A form a fundamental system of entourages of X.

A × A is dense in X × X (Chapter I, § 4, no. 3, Proposition 7). Let V be an open entourage of A; it is the intersection of A × A with an open entourage U of X. We have U ⊂ V̄ (Chapter I, § 1, no. 6, Proposition 5), and this relation together with V̄ ⊂ U establishes the proposition, in view of § 1, no. 2, Corollary 2 of Proposition 2.

